\section{System and Agent Workflow}
\label{sec:system}

\subsection{Services and Agent Operations}
\sys{} combines a web console, a FastAPI backend, a service-management CLI
(\texttt{ragctl}), and an MCP server~\cite{mcp}. The inspected implementation
exposes 41 \texttt{kb\_*} knowledge-management tools within a 94-tool inventory,
including document creation, movement, tagging, search, and offset-addressed
reads. The console calls HTTP services; MCP-capable agents operate the same
resources through tools. These interfaces share data and operations, not
necessarily a reasoning policy: ordinary search endpoints return candidates
without requiring a content-gate judgment.

Figure~\ref{fig:arch} separates the services from the packaged agent skills
that prescribe the demonstrated workflows. ChromaDB stores BGE-M3
vectors~\cite{bge-m3}; Neo4j supports document relations; Markdown and
JSON/YAML metadata retain content and collection structure. Graph navigation
is optional, not the evaluated retrieval route. Deployment requires the
services, model weights, and a supported, configured agent harness.

\subsection{Preparing Addressable Evidence}
MinerU converts PDFs to structured Markdown~\cite{mineru}. Long texts are
split at a configurable character ceiling (30,000 in the supplied corpus).
Retained headings and part identifiers locate a passage within its document;
line-offset reads allow an agent to continue beyond an initial excerpt.

The organization skill asks the agent to inspect parsed content, select among
existing category bases, and use tools to place and tag documents. The
workflow explicitly checks the resulting parts and indexing status:
successful upload alone does not establish complete tagging and indexing.
Administrators retain control over the bases and permissions. The five-base
collection is a demonstration configuration, not a platform restriction.

\subsection{Read, Assess, and Escalate}
QDCVR stands for Query-Driven Content-Verified Retrieval. Its skill rewrites
a query,
recalls vector candidates across selected bases, and reads their text.
\emph{Content-verified} denotes a prompted relevance assessment, not guaranteed
correctness. The executing agent assigns
\[
s=s_{\mathrm{topic}}+s_{\mathrm{scenario}}+s_{\mathrm{evidence}},
\]
with ranges $0$--$3$, $0$--$3$, and $0$--$2$. A score of at least 6 permits
a direct answer. A score-5 candidate is retained as a partial-evidence
backstop; candidates scoring at most 4 are discarded. If none scores at least
6, the librarian inspects base summaries and document metadata, searches
selectively, and reads new candidates. If no stronger evidence emerges, the
skill permits an explicitly attributed answer from a usable score-5 backstop.

Answers contain search paths, answer, sources, confidence, and blind spots;
evidence is addressable by base, document, part, and section. If neither strong evidence
nor a usable backstop remains, the skill requests a not-found report with
scope, near miss, and reason. These instructions rely on agent compliance. The
supplied artifacts retain candidate excerpts and answer records; mutable
file-tree metadata is not an append-only per-query audit log.

